%!TeX root=../emmatop.tex
\chapter[Chapter \thechapter]{}
\lettrine[lraise=0.3]{T}{he} wretchedness of a scheme to Box Hill was in Emma's thoughts all the evening. How it might be considered by the rest of the party, she could not tell. They, in their different homes, and their different ways, might be looking back on it with pleasure; but in her view it was a morning more completely misspent, more totally bare of rational satisfaction at the time, and more to be abhorred in recollection, than any she had ever passed. A whole evening of back-gammon with her father, was felicity to it. There, indeed, lay real pleasure, for there she was giving up the sweetest hours of the twenty-four to his comfort; and feeling that, unmerited as might be the degree of his fond affection and confiding esteem, she could not, in her general conduct, be open to any severe reproach. As a daughter, she hoped she was not without a heart. She hoped no one could have said to her, <How could you be so unfeeling to your father?—I must, I will tell you truths while I can.> Miss Bates should never again—no, never! If attention, in future, could do away the past, she might hope to be forgiven. She had been often remiss, her conscience told her so; remiss, perhaps, more in thought than fact; scornful, ungracious. But it should be so no more. In the warmth of true contrition, she would call upon her the very next morning, and it should be the beginning, on her side, of a regular, equal, kindly intercourse.

She was just as determined when the morrow came, and went early, that nothing might prevent her. It was not unlikely, she thought, that she might see Mr Knightley in her way; or, perhaps, he might come in while she were paying her visit. She had no objection. She would not be ashamed of the appearance of the penitence, so justly and truly hers. Her eyes were towards Donwell as she walked, but she saw him not.

<The ladies were all at home.> She had never rejoiced at the sound before, nor ever before entered the passage, nor walked up the stairs, with any wish of giving pleasure, but in conferring obligation, or of deriving it, except in subsequent ridicule.

There was a bustle on her approach; a good deal of moving and talking. She heard Miss Bates's voice, something was to be done in a hurry; the maid looked frightened and awkward; hoped she would be pleased to wait a moment, and then ushered her in too soon. The aunt and niece seemed both escaping into the adjoining room. Jane she had a distinct glimpse of, looking extremely ill; and, before the door had shut them out, she heard Miss Bates saying, <Well, my dear, I shall say you are laid down upon the bed, and I am sure you are ill enough.>

Poor old Mrs Bates, civil and humble as usual, looked as if she did not quite understand what was going on.

<I am afraid Jane is not very well,> said she, <but I do not know; they tell me she is well. I dare say my daughter will be here presently, Miss Woodhouse. I hope you find a chair. I wish Hetty had not gone. I am very little able—Have you a chair, ma'am? Do you sit where you like? I am sure she will be here presently.>

Emma seriously hoped she would. She had a moment's fear of Miss Bates keeping away from her. But Miss Bates soon came—<Very happy and obliged>—but Emma's conscience told her that there was not the same cheerful volubility as before—less ease of look and manner. A very friendly inquiry after Miss Fairfax, she hoped, might lead the way to a return of old feelings. The touch seemed immediate.

<Ah! Miss Woodhouse, how kind you are!—I suppose you have heard—and are come to give us joy. This does not seem much like joy, indeed, in me—(twinkling away a tear or two)—but it will be very trying for us to part with her, after having had her so long, and she has a dreadful headache just now, writing all the morning:—such long letters, you know, to be written to Colonel Campbell, and Mrs Dixon. <My dear,> said I, <you will blind yourself>—for tears were in her eyes perpetually. One cannot wonder, one cannot wonder. It is a great change; and though she is amazingly fortunate—such a situation, I suppose, as no young woman before ever met with on first going out—do not think us ungrateful, Miss Woodhouse, for such surprising good fortune—(again dispersing her tears)—but, poor dear soul! if you were to see what a headache she has. When one is in great pain, you know one cannot feel any blessing quite as it may deserve. She is as low as possible. To look at her, nobody would think how delighted and happy she is to have secured such a situation. You will excuse her not coming to you—she is not able—she is gone into her own room—I want her to lie down upon the bed. <My dear,> said I, <I shall say you are laid down upon the bed:> but, however, she is not; she is walking about the room. But, now that she has written her letters, she says she shall soon be well. She will be extremely sorry to miss seeing you, Miss Woodhouse, but your kindness will excuse her. You were kept waiting at the door—I was quite ashamed—but somehow there was a little bustle—for it so happened that we had not heard the knock, and till you were on the stairs, we did not know any body was coming. <It is only Mrs Cole,> said I, <depend upon it. Nobody else would come so early.> <Well,> said she, <it must be borne some time or other, and it may as well be now.> But then Patty came in, and said it was you. <Oh!> said I, <it is Miss Woodhouse: I am sure you will like to see her.>—<I can see nobody,> said she; and up she got, and would go away; and that was what made us keep you waiting—and extremely sorry and ashamed we were. <If you must go, my dear,> said I, <you must, and I will say you are laid down upon the bed.>>

Emma was most sincerely interested. Her heart had been long growing kinder towards Jane; and this picture of her present sufferings acted as a cure of every former ungenerous suspicion, and left her nothing but pity; and the remembrance of the less just and less gentle sensations of the past, obliged her to admit that Jane might very naturally resolve on seeing Mrs Cole or any other steady friend, when she might not bear to see herself. She spoke as she felt, with earnest regret and solicitude—sincerely wishing that the circumstances which she collected from Miss Bates to be now actually determined on, might be as much for Miss Fairfax's advantage and comfort as possible. <It must be a severe trial to them all. She had understood it was to be delayed till Colonel Campbell's return.>

<So very kind!> replied Miss Bates. <But you are always kind.>

There was no bearing such an <always;> and to break through her dreadful gratitude, Emma made the direct inquiry of—

<Where—may I ask?—is Miss Fairfax going?>

<To a Mrs Smallridge—charming woman—most superior—to have the charge of her three little girls—delightful children. Impossible that any situation could be more replete with comfort; if we except, perhaps, Mrs Suckling's own family, and Mrs Bragge's; but Mrs Smallridge is intimate with both, and in the very same neighbourhood:—lives only four miles from Maple Grove. Jane will be only four miles from Maple Grove.>

<Mrs Elton, I suppose, has been the person to whom Miss Fairfax owes\longdash>

<Yes, our good Mrs Elton. The most indefatigable, true friend. She would not take a denial. She would not let Jane say, <No;> for when Jane first heard of it, (it was the day before yesterday, the very morning we were at Donwell,) when Jane first heard of it, she was quite decided against accepting the offer, and for the reasons you mention; exactly as you say, she had made up her mind to close with nothing till Colonel Campbell's return, and nothing should induce her to enter into any engagement at present—and so she told Mrs Elton over and over again—and I am sure I had no more idea that she would change her mind!—but that good Mrs Elton, whose judgment never fails her, saw farther than I did. It is not every body that would have stood out in such a kind way as she did, and refuse to take Jane's answer; but she positively declared she would not write any such denial yesterday, as Jane wished her; she would wait—and, sure enough, yesterday evening it was all settled that Jane should go. Quite a surprize to me! I had not the least idea!—Jane took Mrs Elton aside, and told her at once, that upon thinking over the advantages of Mrs Smallridge's situation, she had come to the resolution of accepting it.—I did not know a word of it till it was all settled.>

<You spent the evening with Mrs Elton?>

<Yes, all of us; Mrs Elton would have us come. It was settled so, upon the hill, while we were walking about with Mr Knightley. <You must all spend your evening with us,> said she—<I positively must have you all come.>>

<Mr Knightley was there too, was he?>

<No, not Mr Knightley; he declined it from the first; and though I thought he would come, because Mrs Elton declared she would not let him off, he did not;—but my mother, and Jane, and I, were all there, and a very agreeable evening we had. Such kind friends, you know, Miss Woodhouse, one must always find agreeable, though every body seemed rather fagged after the morning's party. Even pleasure, you know, is fatiguing—and I cannot say that any of them seemed very much to have enjoyed it. However, I shall always think it a very pleasant party, and feel extremely obliged to the kind friends who included me in it.>

<Miss Fairfax, I suppose, though you were not aware of it, had been making up her mind the whole day?>

<I dare say she had.>

<Whenever the time may come, it must be unwelcome to her and all her friends—but I hope her engagement will have every alleviation that is possible—I mean, as to the character and manners of the family.>

<Thank you, dear Miss Woodhouse. Yes, indeed, there is every thing in the world that can make her happy in it. Except the Sucklings and Bragges, there is not such another nursery establishment, so liberal and elegant, in all Mrs Elton's acquaintance. Mrs Smallridge, a most delightful woman!—A style of living almost equal to Maple Grove—and as to the children, except the little Sucklings and little Bragges, there are not such elegant sweet children anywhere. Jane will be treated with such regard and kindness!—It will be nothing but pleasure, a life of pleasure.—And her salary!—I really cannot venture to name her salary to you, Miss Woodhouse. Even you, used as you are to great sums, would hardly believe that so much could be given to a young person like Jane.>

<Ah! madam,> cried Emma, <if other children are at all like what I remember to have been myself, I should think five times the amount of what I have ever yet heard named as a salary on such occasions, dearly earned.>

<You are so noble in your ideas!>

<And when is Miss Fairfax to leave you?>

<Very soon, very soon, indeed; that's the worst of it. Within a fortnight. Mrs Smallridge is in a great hurry. My poor mother does not know how to bear it. So then, I try to put it out of her thoughts, and say, Come ma'am, do not let us think about it any more.>

<Her friends must all be sorry to lose her; and will not Colonel and Mrs Campbell be sorry to find that she has engaged herself before their return?>

<Yes; Jane says she is sure they will; but yet, this is such a situation as she cannot feel herself justified in declining. I was so astonished when she first told me what she had been saying to Mrs Elton, and when Mrs Elton at the same moment came congratulating me upon it! It was before tea—stay—no, it could not be before tea, because we were just going to cards—and yet it was before tea, because I remember thinking—Oh! no, now I recollect, now I have it; something happened before tea, but not that. Mr Elton was called out of the room before tea, old John Abdy's son wanted to speak with him. Poor old John, I have a great regard for him; he was clerk to my poor father twenty-seven years; and now, poor old man, he is bed-ridden, and very poorly with the rheumatic gout in his joints—I must go and see him to-day; and so will Jane, I am sure, if she gets out at all. And poor John's son came to talk to Mr Elton about relief from the parish; he is very well to do himself, you know, being head man at the Crown, ostler, and every thing of that sort, but still he cannot keep his father without some help; and so, when Mr Elton came back, he told us what John ostler had been telling him, and then it came out about the chaise having been sent to Randalls to take Mr Frank Churchill to Richmond. That was what happened before tea. It was after tea that Jane spoke to Mrs Elton.>

Miss Bates would hardly give Emma time to say how perfectly new this circumstance was to her; but as without supposing it possible that she could be ignorant of any of the particulars of Mr Frank Churchill's going, she proceeded to give them all, it was of no consequence.

\begin{figure}[tbph]
\centering
\includegraphics[width=\linewidth]{44chaise}
\caption{Seen the Crown chaise pass by}
\end{figure}

What Mr Elton had learned from the ostler on the subject, being the accumulation of the ostler's own knowledge, and the knowledge of the servants at Randalls, was, that a messenger had come over from Richmond soon after the return of the party from Box Hill—which messenger, however, had been no more than was expected; and that Mr Churchill had sent his nephew a few lines, containing, upon the whole, a tolerable account of Mrs Churchill, and only wishing him not to delay coming back beyond the next morning early; but that Mr Frank Churchill having resolved to go home directly, without waiting at all, and his horse seeming to have got a cold, Tom had been sent off immediately for the Crown chaise, and the ostler had stood out and seen it pass by, the boy going a good pace, and driving very steady.

There was nothing in all this either to astonish or interest, and it caught Emma's attention only as it united with the subject which already engaged her mind. The contrast between Mrs Churchill's importance in the world, and Jane Fairfax's, struck her; one was every thing, the other nothing—and she sat musing on the difference of woman's destiny, and quite unconscious on what her eyes were fixed, till roused by Miss Bates's saying,

<Aye, I see what you are thinking of, the pianoforte. What is to become of that?—Very true. Poor dear Jane was talking of it just now.—<You must go,> said she. <You and I must part. You will have no business here.—Let it stay, however,> said she; <give it houseroom till Colonel Campbell comes back. I shall talk about it to him; he will settle for me; he will help me out of all my difficulties.>—And to this day, I do believe, she knows not whether it was his present or his daughter's.>

Now Emma was obliged to think of the pianoforte; and the remembrance of all her former fanciful and unfair conjectures was so little pleasing, that she soon allowed herself to believe her visit had been long enough; and, with a repetition of every thing that she could venture to say of the good wishes which she really felt, took leave.